\chapter*{总结和展望}
\addcontentsline{toc}{chapter}{总结和展望}

本课程设计围绕“请输入课程设计题目”展开，完成了系统需求分析、总体方案设计、硬件设计、软件设计以及系统测试等工作。通过本次课程设计，基本掌握了相关硬件平台或软件环境的使用方法，并能够根据任务要求完成系统功能划分、模块设计、程序实现和测试验证。

在系统实现方面，本文首先根据课程设计任务书明确了系统的功能需求和设计目标，然后完成了系统总体结构设计，并对主要硬件资源、接口连接关系和软件模块进行了说明。在软件设计过程中，按照模块化设计思路完成了主要功能模块的实现，使系统能够按照预期流程运行。

在系统测试方面，本文对各功能模块和系统整体运行情况进行了测试。测试结果表明，系统能够完成课程设计要求的主要功能，运行现象基本符合预期，说明本文所设计的总体方案和实现方法具有一定可行性。

由于课程设计时间和测试条件有限，系统仍存在一些可以继续改进的地方。例如，可以进一步完善测试数据记录，增加不同环境和不同工况下的重复测试；也可以继续优化程序结构，提高系统运行的稳定性、可维护性和扩展性。后续还可在现有系统基础上增加更多功能模块，使系统具有更完整的应用价值。

综上所述，本次课程设计基本完成了预定任务，达到了课程设计的训练目的。通过本次设计过程，进一步加深了对系统设计、模块划分、软硬件协同调试和实验验证方法的理解，为后续相关课程学习和工程实践打下了基础。